\documentclass[11pt,a4paper]{article}

% ============================================================
% TD PHYSIQUE-CHIMIE TERMINALE -> PREPA
% Version V3 : problemes longs contextualises type concours
% Chimie, thermodynamique, mecanique des fluides
% ============================================================

\usepackage[utf8]{inputenc}
\usepackage[T1]{fontenc}
\usepackage[french]{babel}
\usepackage{lmodern}
\usepackage{geometry}
\geometry{margin=1.45cm}
\usepackage{amsmath,amssymb}
\usepackage{siunitx}
\sisetup{locale=FR,detect-all,per-mode=symbol}
\usepackage{tikz}
\usepackage{tcolorbox}
\tcbuselibrary{skins,breakable}
\usepackage{enumitem}
\usepackage{array,tabularx,booktabs,multicol}
\usepackage{fancyhdr}
\usepackage{hyperref}
\usepackage{xcolor}

\hypersetup{colorlinks=true,linkcolor=blue!50!black,urlcolor=blue!50!black}

\pagestyle{fancy}
\fancyhf{}
\lhead{Terminale vers CPGE - TD concours}
\rhead{Chimie - Thermodynamique - Fluides}
\cfoot{\thepage}

\setlength{\parindent}{0pt}
\setlength{\parskip}{4pt}
\setlist[enumerate]{itemsep=2pt,topsep=3pt}

% ============================================================
% BOITES ET COMMANDES
% ============================================================

\definecolor{bleu}{RGB}{32,76,135}
\definecolor{vert}{RGB}{31,120,72}
\definecolor{rouge}{RGB}{170,45,45}
\definecolor{orange}{RGB}{190,105,25}

\newtcolorbox{methode}{breakable,enhanced,colback=blue!4,colframe=bleu,title=\textbf{Methode},fonttitle=\bfseries}
\newtcolorbox{retenir}{breakable,enhanced,colback=green!4,colframe=vert,title=\textbf{A retenir},fonttitle=\bfseries}
\newtcolorbox{erreur}{breakable,enhanced,colback=red!4,colframe=rouge,title=\textbf{Erreur classique},fonttitle=\bfseries}
\newtcolorbox{vigilance}{breakable,enhanced,colback=orange!5,colframe=orange,title=\textbf{Point de vigilance},fonttitle=\bfseries}
\newtcolorbox{corrige}{breakable,enhanced,colback=gray!5,colframe=black!65,title=\textbf{Corrige detaille},fonttitle=\bfseries}
\newtcolorbox{probleme}[2]{breakable,enhanced,colback=white,colframe=black!75,title=\textbf{Probleme #1 -- #2},fonttitle=\bfseries}

\newcommand{\diffone}{\(\bigstar\star\star\star\star\)}
\newcommand{\difftwo}{\(\bigstar\bigstar\star\star\star\)}
\newcommand{\diffthree}{\(\bigstar\bigstar\bigstar\star\star\)}
\newcommand{\difffour}{\(\bigstar\bigstar\bigstar\bigstar\star\)}
\newcommand{\difffive}{\(\bigstar\bigstar\bigstar\bigstar\bigstar\)}

\newcommand{\pH}{\mathrm{pH}}
\newcommand{\pKa}{\mathrm{p}K_a}
\newcommand{\e}{\mathrm{e}}

\begin{document}

% ============================================================
% PAGE DE TITRE
% ============================================================

\begin{titlepage}
\centering
\vspace*{0.8cm}
{\Huge \bfseries TD concours de physique-chimie\par}
\vspace{0.25cm}
{\LARGE Terminale tres exigeant - preparation a la CPGE\par}
\vspace{0.55cm}
{\Large \bfseries Chimie, thermodynamique et mecanique des fluides\par}
\vspace{0.35cm}
{\large Version longue : problemes contextualises, transversaux, corriges detailles\par}

\vspace{0.8cm}

\begin{tcolorbox}[colback=blue!4,colframe=bleu,width=0.93\textwidth]
\textbf{Objectif.} Ce document transforme les exercices courts en \textbf{problemes longs de type concours}. Chaque probleme suit un fil directeur : contexte scientifique ou technique, donnees realistes, plusieurs parties, bilan de matiere, bilan d'energie, pression, debit, rendement, critique de modele et interpretation physique.

\medskip

\textbf{Niveau.} Fin de terminale tres solide, eleves visant une classe preparatoire. Les notions de fond restent celles de terminale : quantite de matiere, titrage, transformation chimique, energie thermique, puissance, rendement, hydrostatique, debit, Bernoulli simplifie et poussee d'Archimede. La difficulte vient de la longueur, des enchainements et du sens physique.

\medskip

\textbf{Gradient.} Les difficultes sont indiquees par etoiles. Les problemes \difffour{} et \difffive{} sont volontairement proches de l'esprit CPGE : modelisation, ordres de grandeur, hypotheses et critiques.
\end{tcolorbox}

\vfill
\textbf{Conseil eleve.} Ne pas chercher une formule magique. Lire, schematiser, convertir, poser les bilans, calculer, interpreter.
\end{titlepage}

\tableofcontents
\newpage

% ============================================================
% RAPPELS
% ============================================================

\section{Rappels tres courts de methode}

\begin{retenir}
\textbf{Chimie.}
\[
 n=\frac{m}{M},\qquad C=\frac{n}{V},\qquad C_m=\frac{m}{V},\qquad \sigma=\sum_i \lambda_i [X_i].
\]
Pour un titrage a l'equivalence :
\[
\frac{n_{\text{titre}}}{\nu_{\text{titre}}}=\frac{n_{\text{titrant}}}{\nu_{\text{titrant}}}.
\]
Pour un acide faible :
\[
K_a=\frac{[A^-][H_3O^+]}{[HA]},\quad \pH=-\log[H_3O^+],\quad \pKa=-\log K_a.
\]
\end{retenir}

\begin{retenir}
\textbf{Thermodynamique et energie.}
\[
Q=mc\Delta T,\qquad Q=mL,\qquad E=n|\Delta_rH|,\qquad \eta=\frac{E_{\text{utile}}}{E_{\text{fournie}}}.
\]
Si le systeme est thermiquement isole :
\[
\sum_i Q_i=0.
\]
\end{retenir}

\begin{retenir}
\textbf{Fluides.}
\[
p=p_0+\rho gh,\qquad F=pS,\qquad F_A=\rho_{\text{fluide}}gV_{\text{immerge}},\qquad D_v=Sv.
\]
Bernoulli simplifie :
\[
p+\frac12\rho v^2+\rho gz=\text{constante}.
\]
\end{retenir}

\begin{methode}
\textbf{Strategie pour les problemes longs.}
\begin{enumerate}
\item Identifier le systeme : solution, fluide, reservoir, gaz, eau chauffee, objet immerge.
\item Convertir toutes les donnees avant de calculer.
\item Chercher les bilans : matiere, energie, pression, debit.
\item Garder les expressions litterales jusqu'a la fin.
\item Verifier l'unite et l'ordre de grandeur.
\item Terminer par une phrase d'interpretation et une critique des hypotheses.
\end{enumerate}
\end{methode}

\begin{erreur}
\textbf{Signaux d'alarme.} Rendement superieur a 1, temperature finale hors intervalle dans un melange, pression negative, force pressante ridiculement faible a grande profondeur, debit enorme dans un petit tuyau, oubli d'un facteur de dilution ou d'un coefficient stoechiometrique.
\end{erreur}

\newpage

% ============================================================
% ENONCES
% ============================================================

\section{Enonces : problemes longs contextualises}

\subsection{Niveau \diffone{} - problemes longs de prise en main}

\begin{probleme}{1 \diffone}{Controle d'une solution saline pour conductimetrie}
Un laboratoire de lycee prepare une gamme de solutions de chlorure de sodium pour verifier le fonctionnement d'un conductimetre. La solution mere doit avoir un volume \(V_m=\SI{250}{mL}\) et une concentration \(C_m=\SI{0.100}{mol.L^{-1}}\). On prepare ensuite une solution fille de volume \(V_f=\SI{100.0}{mL}\) et de concentration \(C_f=\SI{1.00e-2}{mol.L^{-1}}\). La conductivite mesuree de la solution fille vaut \(\sigma_{\text{mes}}=\SI{0.119}{S.m^{-1}}\).

Donnees : \(M(NaCl)=\SI{58.5}{g.mol^{-1}}\), \(\lambda_{Na^+}=\SI{5.01e-3}{S.m^2.mol^{-1}}\), \(\lambda_{Cl^-}=\SI{7.63e-3}{S.m^2.mol^{-1}}\).

\textbf{Partie A - Preparation.}
\begin{enumerate}
\item Calculer la quantite de matiere de \(NaCl\) necessaire pour la solution mere.
\item Calculer la masse de solide a peser.
\item Rediger un protocole precis en mentionnant la verrerie.
\end{enumerate}

\textbf{Partie B - Dilution.}
\begin{enumerate}[resume]
\item Determiner le facteur de dilution.
\item Calculer le volume de solution mere a prelever.
\item Expliquer pourquoi une pipette jaugee est preferable a une eprouvette graduee.
\end{enumerate}

\textbf{Partie C - Conductivite.}
\begin{enumerate}[resume]
\item Convertir \(C_f\) en \(\si{mol.m^{-3}}\).
\item Calculer la conductivite theorique \(\sigma_{\text{th}}\).
\item Calculer l'ecart relatif entre mesure et theorie.
\item Proposer deux causes experimentales expliquant l'ecart.
\end{enumerate}
\end{probleme}

\begin{probleme}{2 \diffone}{Bouilloire, rendement et extrapolation domestique}
Une bouilloire de puissance nominale \(P=\SI{1800}{W}\) chauffe \(m=\SI{0.750}{kg}\) d'eau de \(\SI{18.0}{\celsius}\) a \(\SI{95.0}{\celsius}\) en \(\SI{160}{s}\). On veut ensuite estimer le temps necessaire pour chauffer \(\SI{1.20}{kg}\) d'eau de \(\SI{20.0}{\celsius}\) a \(\SI{100.0}{\celsius}\).

Donnee : \(c_{\text{eau}}=\SI{4180}{J.kg^{-1}.K^{-1}}\).

\textbf{Partie A - Premiere experience.}
\begin{enumerate}
\item Calculer l'energie thermique recue par l'eau.
\item Calculer l'energie electrique consommee.
\item Determiner le rendement de la bouilloire.
\end{enumerate}

\textbf{Partie B - Prediction.}
\begin{enumerate}[resume]
\item Calculer l'energie utile necessaire pour la seconde chauffe.
\item En supposant le meme rendement, calculer l'energie electrique a fournir.
\item En deduire la duree de chauffage.
\end{enumerate}

\textbf{Partie C - Discussion.}
\begin{enumerate}[resume]
\item Pourquoi le rendement peut-il changer avec la masse d'eau ?
\item Expliquer pourquoi la temperature d'ebullition n'est pas toujours exactement \(\SI{100}{\celsius}\).
\item Citer deux pertes thermiques possibles.
\end{enumerate}
\end{probleme}

\begin{probleme}{3 \diffone}{Piscine municipale : pression, force et securite}
Dans une piscine, une trappe circulaire de rayon \(r=\SI{12.0}{cm}\) est situee au fond d'une zone de profondeur \(h=\SI{3.20}{m}\). On veut estimer la force exercee par l'eau sur cette trappe et comprendre pourquoi les dispositifs de vidange doivent etre securises.

Donnees : \(\rho=\SI{1000}{kg.m^{-3}}\), \(g=\SI{9.81}{m.s^{-2}}\), \(p_{\text{atm}}=\SI{1.01e5}{Pa}\).

\textbf{Partie A - Pression.}
\begin{enumerate}
\item Calculer la surpression au fond.
\item Calculer la pression absolue au fond.
\item Convertir en bar.
\end{enumerate}

\textbf{Partie B - Force pressante.}
\begin{enumerate}[resume]
\item Calculer l'aire de la trappe.
\item Calculer la force pressante absolue.
\item Calculer la force due uniquement a la surpression de l'eau.
\item Comparer cette derniere au poids d'une personne de masse \(\SI{70}{kg}\).
\end{enumerate}

\textbf{Partie C - Interpretation.}
\begin{enumerate}[resume]
\item Pourquoi utilise-t-on souvent la surpression pour raisonner sur les efforts mecaniques ?
\item Pourquoi la pression ne depend-elle pas de la largeur de la piscine ?
\end{enumerate}
\end{probleme}

\begin{probleme}{4 \diffone}{Détartrant menager : dilution, titrage et concentration massique}
Un detartrant contient de l'acide chlorhydrique. On preleve \(\SI{5.00}{mL}\) de produit commercial et on dilue a \(\SI{100.0}{mL}\). On titre ensuite \(\SI{10.0}{mL}\) de solution diluee par une solution de soude de concentration \(C_B=\SI{0.100}{mol.L^{-1}}\). L'equivalence est atteinte pour \(V_E=\SI{14.8}{mL}\).

Donnee : \(M(HCl)=\SI{36.5}{g.mol^{-1}}\).

\textbf{Partie A - Reaction de titrage.}
\begin{enumerate}
\item Ecrire l'equation de reaction entre \(H_3O^+\) et \(HO^-\).
\item Expliquer le sens physique de l'equivalence.
\end{enumerate}

\textbf{Partie B - Exploitation.}
\begin{enumerate}[resume]
\item Calculer la quantite de matiere de soude versee a l'equivalence.
\item En deduire la concentration de la solution diluee.
\item Determiner le facteur de dilution.
\item Calculer la concentration molaire du produit commercial.
\end{enumerate}

\textbf{Partie C - Controle qualite.}
\begin{enumerate}[resume]
\item Calculer la concentration massique en \(HCl\).
\item Un fabricant annonce \(\SI{110}{g.L^{-1}}\). Comparer.
\item Donner deux precautions de securite lors de la manipulation.
\end{enumerate}
\end{probleme}

\begin{probleme}{5 \diffone}{Fontaine publique : debit, vitesse et bilan journalier}
Une fontaine publique remplit une bouteille de \(\SI{1.50}{L}\) en \(\SI{12.0}{s}\). L'eau sort par un tube cylindrique de diametre interieur \(d=\SI{8.0}{mm}\). La ville souhaite estimer le volume d'eau distribue si la fontaine fonctionne \(\SI{3.0}{h}\) par jour.

\textbf{Partie A - Debit.}
\begin{enumerate}
\item Calculer le debit volumique en \(\si{L.s^{-1}}\).
\item Convertir en \(\si{m^3.s^{-1}}\).
\end{enumerate}

\textbf{Partie B - Vitesse.}
\begin{enumerate}[resume]
\item Calculer la section du tube.
\item En deduire la vitesse moyenne de l'eau.
\item Commenter l'ordre de grandeur.
\end{enumerate}

\textbf{Partie C - Usage municipal.}
\begin{enumerate}[resume]
\item Calculer le volume distribue en une journee.
\item Calculer le volume distribue en \(30\) jours.
\item Pourquoi ce calcul surestime-t-il probablement la consommation reelle ?
\end{enumerate}
\end{probleme}

\begin{probleme}{6 \diffone}{Glaçons dans une boisson : bilan thermique complet}
On place \(m_g=\SI{40}{g}\) de glace a \(\SI{0}{\celsius}\) dans \(m_e=\SI{300}{g}\) d'eau a \(\SI{22}{\celsius}\). Le recipient est suppose parfaitement isolant.

Donnees : \(L_f=\SI{334}{kJ.kg^{-1}}\), \(c_{\text{eau}}=\SI{4180}{J.kg^{-1}.K^{-1}}\).

\textbf{Partie A - Possibilite de fusion totale.}
\begin{enumerate}
\item Calculer l'energie necessaire pour fondre toute la glace.
\item Calculer l'energie maximale que l'eau liquide peut ceder en descendant a \(\SI{0}{\celsius}\).
\item Conclure sur la fusion totale.
\end{enumerate}

\textbf{Partie B - Temperature finale.}
\begin{enumerate}[resume]
\item Calculer l'energie restante apres fusion.
\item Calculer la temperature finale.
\item Verifier que le resultat est coherent.
\end{enumerate}

\textbf{Partie C - Variante.}
\begin{enumerate}[resume]
\item Que se passerait-il si la masse de glace etait doublee ?
\item Calculer alors la masse de glace restante a \(\SI{0}{\celsius}\).
\end{enumerate}
\end{probleme}

\begin{probleme}{7 \diffone}{Ballon d'helium : poussee d'Archimede et charge utile}
Un ballon d'helium de volume \(V=\SI{0.030}{m^3}\) possede une enveloppe de masse \(m_{env}=\SI{8.0}{g}\). Il doit soulever une petite carte publicitaire.

Donnees : \(\rho_{air}=\SI{1.20}{kg.m^{-3}}\), \(\rho_{He}=\SI{0.18}{kg.m^{-3}}\), \(g=\SI{9.81}{m.s^{-2}}\).

\textbf{Partie A - Forces.}
\begin{enumerate}
\item Calculer la poussee d'Archimede.
\item Calculer la masse d'helium.
\item Calculer le poids ballon + helium.
\end{enumerate}

\textbf{Partie B - Decollage.}
\begin{enumerate}[resume]
\item Le ballon decolle-t-il sans charge ?
\item Calculer la masse maximale de charge pouvant etre soulevee.
\item Pourquoi cette valeur est-elle une valeur maximale ideale ?
\end{enumerate}

\textbf{Partie C - Sens physique.}
\begin{enumerate}[resume]
\item Que se passe-t-il si le ballon fuit lentement ?
\item Que se passe-t-il si l'air exterieur devient moins dense ?
\end{enumerate}
\end{probleme}

\begin{probleme}{8 \diffone}{Rechaud au propane : chimie, energie et dioxyde de carbone}
Un rechaud de camping brule \(m=\SI{5.00}{g}\) de propane \(C_3H_8\). La combustion complete libere \(\SI{2220}{kJ.mol^{-1}}\). Le rendement de transfert a une casserole d'eau est \(\eta=45\%\).

Donnees : \(M(C_3H_8)=\SI{44.0}{g.mol^{-1}}\), \(M(CO_2)=\SI{44.0}{g.mol^{-1}}\), \(c_{eau}=\SI{4180}{J.kg^{-1}.K^{-1}}\).

\textbf{Partie A - Reaction.}
\begin{enumerate}
\item Ajuster la combustion complete du propane.
\item Calculer la quantite de propane brulee.
\item Calculer la masse de \(CO_2\) produite.
\end{enumerate}

\textbf{Partie B - Chauffage.}
\begin{enumerate}[resume]
\item Calculer l'energie chimique liberee.
\item Calculer l'energie utile recue par l'eau.
\item Determiner la masse d'eau chauffable de \(\SI{20}{\celsius}\) a \(\SI{80}{\celsius}\).
\end{enumerate}

\textbf{Partie C - Analyse.}
\begin{enumerate}[resume]
\item Pourquoi le rendement n'est-il pas de \(100\%\) ?
\item Pourquoi est-il dangereux d'utiliser ce rechaud dans une tente fermee ?
\end{enumerate}
\end{probleme}

\subsection{Niveau \difftwo{} - problemes standard approfondis}

\begin{probleme}{9 \difftwo}{Controle d'un vinaigre : dilution, titrage, incertitudes}
Un vinaigre contient de l'acide ethanoique. On preleve \(\SI{10.0}{mL}\) de vinaigre et on dilue a \(\SI{100.0}{mL}\). On titre \(\SI{20.0}{mL}\) de solution diluee par une solution de soude de concentration \(\SI{0.100}{mol.L^{-1}}\). L'equivalence est obtenue pour \(\SI{16.4}{mL}\).

\[
CH_3COOH+HO^-\rightarrow CH_3COO^-+H_2O.
\]
Donnee : \(M(CH_3COOH)=\SI{60.0}{g.mol^{-1}}\).

\textbf{Partie A - Exploitation du dosage.}
\begin{enumerate}
\item Justifier la relation a l'equivalence.
\item Calculer la concentration molaire de la solution diluee.
\item En deduire celle du vinaigre.
\end{enumerate}

\textbf{Partie B - Degré du vinaigre.}
\begin{enumerate}[resume]
\item Calculer la concentration massique.
\item Comparer a un vinaigre annonce a \(\SI{50}{g.L^{-1}}\).
\item Calculer l'ecart relatif.
\end{enumerate}

\textbf{Partie C - Incertitudes.}
\begin{enumerate}[resume]
\item Identifier trois sources d'incertitude.
\item Expliquer pourquoi la dilution est utile experimentalement.
\item Que se passe-t-il si on depasse legerement l'equivalence ?
\end{enumerate}
\end{probleme}

\begin{probleme}{10 \difftwo}{Chauffe-eau solaire : stockage, rendement et melange}
Un ballon solaire contient \(\SI{150}{L}\) d'eau. Une journee ensoleillee fait passer sa temperature de \(\SI{18}{\celsius}\) a \(\SI{52}{\celsius}\). L'energie solaire incidente sur les capteurs vaut \(\SI{26}{MJ}\). Une douche utilise \(\SI{45}{L}\) d'eau a \(\SI{40}{\celsius}\), obtenue par melange d'eau chaude a \(\SI{52}{\celsius}\) et d'eau froide a \(\SI{15}{\celsius}\).

Donnees : \(\rho=\SI{1000}{kg.m^{-3}}\), \(c=\SI{4180}{J.kg^{-1}.K^{-1}}\).

\textbf{Partie A - Energie stockee.}
\begin{enumerate}
\item Calculer la masse d'eau du ballon.
\item Calculer l'energie thermique stockee.
\item Exprimer cette energie en \(\si{kWh}\).
\end{enumerate}

\textbf{Partie B - Rendement.}
\begin{enumerate}[resume]
\item Determiner le rendement global.
\item Commenter la valeur obtenue.
\end{enumerate}

\textbf{Partie C - Melange.}
\begin{enumerate}[resume]
\item Poser un bilan thermique pour la douche.
\item Calculer le volume d'eau chaude necessaire.
\item Estimer le nombre maximal de douches.
\end{enumerate}
\end{probleme}

\begin{probleme}{11 \difftwo}{Perfusion gravitaire : hauteur, debit et pertes cachees}
Une perfusion necessite une surpression de \(\SI{3.0e3}{Pa}\) a l'entree d'un catheter. La solution est assimilee a de l'eau salee de masse volumique \(\rho=\SI{1030}{kg.m^{-3}}\). Le debit prescrit est \(\SI{2.5}{mL.min^{-1}}\) dans un tube de section \(S=\SI{0.75}{mm^2}\).

\textbf{Partie A - Hydrostatique.}
\begin{enumerate}
\item Etablir la hauteur minimale de la poche.
\item Calculer cette hauteur.
\item Comparer a une installation hospitaliere courante.
\end{enumerate}

\textbf{Partie B - Cinematique du fluide.}
\begin{enumerate}[resume]
\item Convertir le debit en \(\si{m^3.s^{-1}}\).
\item Calculer la vitesse moyenne dans le tube.
\item Calculer la duree necessaire pour vider une poche de \(\SI{500}{mL}\).
\end{enumerate}

\textbf{Partie C - Analyse critique.}
\begin{enumerate}[resume]
\item Pourquoi la hauteur ne suffit-elle pas a determiner le debit ?
\item Citer deux parametres geometriques du tube qui influencent les pertes.
\item Que se passe-t-il si le patient leve le bras ?
\end{enumerate}
\end{probleme}

\begin{probleme}{12 \difftwo}{Comprime anti-calcaire : reaction, gaz et exces d'acide}
Un comprime anti-calcaire de masse \(\SI{1.80}{g}\) contient \(75\%\) de carbonate de calcium. On le plonge dans \(\SI{100}{mL}\) d'acide chlorhydrique de concentration \(\SI{0.500}{mol.L^{-1}}\).

\[
CaCO_3+2H_3O^+\rightarrow Ca^{2+}+CO_2+3H_2O.
\]
Donnees : \(M(CaCO_3)=\SI{100}{g.mol^{-1}}\), \(M(CO_2)=\SI{44.0}{g.mol^{-1}}\).

\textbf{Partie A - Composition du comprime.}
\begin{enumerate}
\item Calculer la masse de carbonate de calcium.
\item Calculer la quantite de carbonate.
\item Calculer la quantite initiale d'ions oxonium.
\end{enumerate}

\textbf{Partie B - Avancement.}
\begin{enumerate}[resume]
\item Identifier le reactif limitant.
\item Calculer la quantite de \(CO_2\) formee.
\item Calculer la masse de \(CO_2\) formee.
\end{enumerate}

\textbf{Partie C - Gaz.}
\begin{enumerate}[resume]
\item En prenant \(V_m=\SI{24.0}{L.mol^{-1}}\), calculer le volume de gaz.
\item Pourquoi observe-t-on une effervescence ?
\item Que changerait un acide deux fois moins concentre ?
\end{enumerate}
\end{probleme}

\begin{probleme}{13 \difftwo}{Refroidissement d'un ordinateur portable}
Un ordinateur portable dissipe une puissance thermique moyenne de \(\SI{45}{W}\). Un ventilateur fait circuler de l'air entrant a \(\SI{22}{\celsius}\) et sortant a \(\SI{38}{\celsius}\). La surface utile de la grille est \(\SI{12}{cm^2}\).

Donnees : \(\rho_{air}=\SI{1.20}{kg.m^{-3}}\), \(c_{air}=\SI{1000}{J.kg^{-1}.K^{-1}}\).

\textbf{Partie A - Debit thermique.}
\begin{enumerate}
\item Etablir la relation \(P=\dot m c\Delta T\).
\item Calculer le debit massique d'air.
\item Calculer le debit volumique.
\end{enumerate}

\textbf{Partie B - Vitesse.}
\begin{enumerate}[resume]
\item Convertir la surface de grille.
\item Calculer la vitesse moyenne de l'air.
\item Commenter le bruit possible du ventilateur.
\end{enumerate}

\textbf{Partie C - Modele.}
\begin{enumerate}[resume]
\item Pourquoi la temperature de sortie n'est-elle pas uniforme ?
\item Quels modes de transfert thermique sont oublies ?
\item Pourquoi la poussiere degrade-t-elle le refroidissement ?
\end{enumerate}
\end{probleme}

\begin{probleme}{14 \difftwo}{Bouteille de plongee : gaz parfait, pression et autonomie}
Une bouteille de plongee de volume \(\SI{12.0}{L}\) contient de l'air a \(\SI{200}{bar}\) et \(T=\SI{293}{K}\). On assimile l'air a un gaz parfait. Le plongeur consomme \(\SI{20}{L.min^{-1}}\) d'air ramene a \(\SI{1.0}{bar}\).

Donnees : \(R=\SI{8.314}{J.mol^{-1}.K^{-1}}\), \(p_{atm}=\SI{1.0e5}{Pa}\).

\textbf{Partie A - Quantite d'air.}
\begin{enumerate}
\item Convertir les donnees en unites SI.
\item Calculer la quantite de matiere dans la bouteille.
\item Calculer le volume equivalent a \(\SI{1.0}{bar}\).
\end{enumerate}

\textbf{Partie B - Autonomie.}
\begin{enumerate}[resume]
\item Estimer l'autonomie en surface.
\item Estimer la pression absolue a \(\SI{20}{m}\) de profondeur.
\item Expliquer pourquoi l'autonomie diminue en profondeur.
\end{enumerate}

\textbf{Partie C - Securite.}
\begin{enumerate}[resume]
\item Pourquoi ne doit-on jamais vider completement une bouteille ?
\item Pourquoi la remontee doit-elle etre controlee ?
\end{enumerate}
\end{probleme}

\begin{probleme}{15 \difftwo}{Neutralisation d'un rejet acide : chimie et echauffement}
Une cuve contient \(\SI{80}{L}\) d'une solution acide de concentration \(C_A=\SI{2.0e-2}{mol.L^{-1}}\) en ions \(H_3O^+\). On veut la neutraliser avec une solution de soude de concentration \(C_B=\SI{0.50}{mol.L^{-1}}\). La neutralisation libere environ \(\SI{57}{kJ.mol^{-1}}\).

\textbf{Partie A - Volume de soude.}
\begin{enumerate}
\item Ecrire la reaction de neutralisation.
\item Calculer la quantite d'acide dans la cuve.
\item Calculer le volume de soude a verser.
\end{enumerate}

\textbf{Partie B - Energie.}
\begin{enumerate}[resume]
\item Calculer l'energie liberee.
\item Estimer l'elevation de temperature si la solution a une masse de \(\SI{80}{kg}\).
\item Donnee : \(c\approx\SI{4180}{J.kg^{-1}.K^{-1}}\).
\end{enumerate}

\textbf{Partie C - Environnement.}
\begin{enumerate}[resume]
\item Pourquoi un exces de soude est-il problematique ?
\item Pourquoi mesure-t-on le pH apres melange ?
\item Quelles precautions experimentales prendre ?
\end{enumerate}
\end{probleme}

\begin{probleme}{16 \difftwo}{Citerne surelevee et arrosage gravitaire}
Une citerne est placee sur une tour. Le niveau libre de l'eau est a \(\SI{6.0}{m}\) au-dessus de l'embout d'arrosage, de diametre \(\SI{10}{mm}\). On neglige d'abord les pertes.

\textbf{Partie A - Pression disponible.}
\begin{enumerate}
\item Calculer la surpression a l'embout.
\item Convertir en bar.
\item Commenter la faiblesse ou non de cette pression.
\end{enumerate}

\textbf{Partie B - Bernoulli.}
\begin{enumerate}[resume]
\item Enoncer les hypotheses permettant d'ecrire \(v=\sqrt{2gh}\).
\item Calculer la vitesse de sortie.
\item Calculer le debit en \(\si{L.min^{-1}}\).
\end{enumerate}

\textbf{Partie C - Realite.}
\begin{enumerate}[resume]
\item Pourquoi le debit reel est-il plus faible ?
\item Quel est l'effet d'un tuyau plus long ?
\item Quel est l'effet d'un embout plus fin ?
\end{enumerate}
\end{probleme}

\subsection{Niveau \diffthree{} - problemes transversaux}

\begin{probleme}{17 \diffthree}{Potabilisation d'urgence : dilution, dosage et controle}
Une equipe humanitaire doit traiter \(\SI{500}{L}\) d'eau avec une solution d'hypochlorite de sodium. La solution commerciale a une concentration \(C_0=\SI{0.50}{mol.L^{-1}}\). La concentration visee dans la cuve est \(C_f=\SI{2.0e-4}{mol.L^{-1}}\). Pour verifier le traitement, on dose un prelevement de \(\SI{50.0}{mL}\) et on trouve \(n=\SI{9.0e-6}{mol}\) d'ions hypochlorite.

\textbf{Partie A - Dose directe.}
\begin{enumerate}
\item Calculer la quantite de matiere necessaire dans la cuve.
\item Calculer le volume de solution commerciale a ajouter.
\item Commenter la precision necessaire.
\end{enumerate}

\textbf{Partie B - Solution intermediaire.}
\begin{enumerate}[resume]
\item On dilue dix fois la solution commerciale. Calculer la nouvelle concentration.
\item Calculer le volume de solution intermediaire a utiliser.
\item Expliquer pourquoi cette methode est plus robuste sur le terrain.
\end{enumerate}

\textbf{Partie C - Controle.}
\begin{enumerate}[resume]
\item Calculer la concentration mesuree.
\item Calculer l'ecart relatif par rapport a la cible.
\item Proposer deux causes physico-chimiques d'ecart.
\end{enumerate}
\end{probleme}

\begin{probleme}{18 \diffthree}{Bain-marie : chauffage, recipient et maintien en temperature}
Un bain-marie contient \(\SI{2.00}{kg}\) d'eau. Une resistance de puissance \(\SI{600}{W}\) chauffe l'eau de \(\SI{20.0}{\celsius}\) a \(\SI{60.0}{\celsius}\) en \(\SI{12.0}{min}\). Le recipient a une capacite thermique \(C_{rec}=\SI{650}{J.K^{-1}}\). Une fois a \(\SI{60}{\celsius}\), la resistance fournit \(\SI{90}{W}\) pour maintenir la temperature.

Donnee : \(c=\SI{4180}{J.kg^{-1}.K^{-1}}\).

\textbf{Partie A - Rendement apparent.}
\begin{enumerate}
\item Calculer l'energie recue par l'eau.
\item Calculer l'energie electrique fournie.
\item Determiner le rendement apparent.
\end{enumerate}

\textbf{Partie B - Correction par le recipient.}
\begin{enumerate}[resume]
\item Calculer l'energie recue par le recipient.
\item Recalculer le rendement utile eau + recipient.
\item Interpreter la difference.
\end{enumerate}

\textbf{Partie C - Maintien.}
\begin{enumerate}[resume]
\item Que represente la puissance \(\SI{90}{W}\) ?
\item Calculer l'energie perdue en \(\SI{30}{min}\).
\item Pourquoi les pertes augmentent-elles quand la temperature du bain augmente ?
\end{enumerate}
\end{probleme}

\begin{probleme}{19 \diffthree}{Fermentation : stoechiometrie, gaz et chaleur}
Dans un fermenteur, du glucose se transforme selon :
\[
C_6H_{12}O_6\rightarrow2C_2H_6O+2CO_2.
\]
Le fermenteur contient initialement \(\SI{1.50}{kg}\) de glucose. On suppose que \(65\%\) du glucose est consomme. La reaction libere environ \(\SI{70}{kJ}\) par mole de glucose consommee. Le moût a une masse \(\SI{20}{kg}\) et une capacite thermique massique \(\SI{4000}{J.kg^{-1}.K^{-1}}\).

Donnees : \(M(C_6H_{12}O_6)=\SI{180}{g.mol^{-1}}\), \(M(CO_2)=\SI{44.0}{g.mol^{-1}}\), \(V_m=\SI{24.0}{L.mol^{-1}}\).

\textbf{Partie A - Chimie.}
\begin{enumerate}
\item Calculer la quantite initiale de glucose.
\item Calculer la quantite consommee.
\item Calculer la quantite de \(CO_2\) produite.
\end{enumerate}

\textbf{Partie B - Gaz.}
\begin{enumerate}[resume]
\item Calculer la masse de \(CO_2\) produite.
\item Calculer le volume de gaz degage.
\item Pourquoi le fermenteur doit-il etre ventile ?
\end{enumerate}

\textbf{Partie C - Thermique.}
\begin{enumerate}[resume]
\item Calculer l'energie thermique liberee.
\item Estimer l'elevation de temperature du moût sans refroidissement.
\item Pourquoi un controle thermique est-il necessaire en industrie ?
\end{enumerate}
\end{probleme}

\begin{probleme}{20 \diffthree}{Chauffe-eau instantane : debit, puissance et faisabilite}
Un chauffe-eau instantane doit fournir \(\SI{6.0}{L.min^{-1}}\) d'eau a \(\SI{42}{\celsius}\), a partir d'eau froide a \(\SI{12}{\celsius}\). Le rendement est \(92\%\). Une prise domestique standard fournit au maximum \(\SI{3.7}{kW}\).

Donnees : \(\rho=\SI{1000}{kg.m^{-3}}\), \(c=\SI{4180}{J.kg^{-1}.K^{-1}}\).

\textbf{Partie A - Puissance thermique.}
\begin{enumerate}
\item Convertir le debit en \(\si{m^3.s^{-1}}\).
\item Calculer le debit massique.
\item Calculer la puissance thermique utile.
\end{enumerate}

\textbf{Partie B - Puissance electrique.}
\begin{enumerate}[resume]
\item Calculer la puissance electrique requise.
\item Comparer a la puissance disponible.
\item Conclure.
\end{enumerate}

\textbf{Partie C - Optimisation.}
\begin{enumerate}[resume]
\item A puissance disponible \(\SI{3.7}{kW}\), quel debit maximal permettrait la meme elevation de temperature ?
\item Pourquoi les chauffe-eau instantanes puissants demandent-ils une installation specifique ?
\end{enumerate}
\end{probleme}

\begin{probleme}{21 \diffthree}{Sous-marin miniature : ballast, pression et hublot}
Un sous-marin miniature de volume exterieur \(V=\SI{0.080}{m^3}\) a une masse ballast vide \(m_0=\SI{70.0}{kg}\). Il peut embarquer de l'eau dans un ballast de volume maximal \(\SI{12.0}{L}\). Il possede un hublot circulaire de rayon \(\SI{8.0}{cm}\).

Donnees : \(\rho=\SI{1000}{kg.m^{-3}}\), \(g=\SI{9.81}{m.s^{-2}}\), \(p_{atm}=\SI{1.01e5}{Pa}\).

\textbf{Partie A - Flottabilite.}
\begin{enumerate}
\item Calculer le poids ballast vide.
\item Calculer la poussee d'Archimede en immersion totale.
\item Le sous-marin remonte-t-il ou descend-il ?
\end{enumerate}

\textbf{Partie B - Equilibre.}
\begin{enumerate}[resume]
\item Determiner la masse d'eau a embarquer pour etre a l'equilibre.
\item Le ballast maximal suffit-il ?
\item Comment piloter une profondeur stable ?
\end{enumerate}

\textbf{Partie C - Pression.}
\begin{enumerate}[resume]
\item Calculer la pression absolue a \(\SI{25}{m}\).
\item Calculer la force sur le hublot.
\item Pourquoi la conception mecanique devient-elle critique ?
\end{enumerate}
\end{probleme}

\begin{probleme}{22 \diffthree}{Moteur thermique : carburant, puissance et refroidissement}
Un petit moteur consomme \(\SI{0.80}{L}\) d'essence par heure. On assimile l'essence a de l'octane \(C_8H_{18}\) de masse volumique \(\SI{750}{kg.m^{-3}}\). La combustion d'une mole d'octane libere \(\SI{5470}{kJ}\). Le moteur convertit \(25\%\) de l'energie chimique en energie mecanique utile, le reste etant evacue sous forme thermique. L'eau de refroidissement entre a \(\SI{20}{\celsius}\) et sort a \(\SI{70}{\celsius}\).

Donnee : \(M(C_8H_{18})=\SI{114}{g.mol^{-1}}\).

\textbf{Partie A - Carburant.}
\begin{enumerate}
\item Ajuster l'equation de combustion complete de l'octane.
\item Calculer la masse de carburant consommee par heure.
\item Calculer la quantite de matiere consommee par heure.
\end{enumerate}

\textbf{Partie B - Puissances.}
\begin{enumerate}[resume]
\item Calculer l'energie chimique fournie par heure.
\item Calculer la puissance chimique moyenne.
\item Calculer la puissance mecanique utile.
\item Calculer la puissance thermique a evacuer.
\end{enumerate}

\textbf{Partie C - Refroidissement.}
\begin{enumerate}[resume]
\item Calculer le debit massique d'eau necessaire.
\item Convertir en \(\si{L.min^{-1}}\).
\item Pourquoi un vrai circuit utilise-t-il un radiateur et un thermostat ?
\end{enumerate}
\end{probleme}

\begin{probleme}{23 \diffthree}{Barrage : debit, energie et conduite forcee}
Un barrage alimente une turbine avec un debit \(D_v=\SI{0.80}{m^3.s^{-1}}\) et une hauteur de chute \(h=\SI{12}{m}\). Le rendement global turbine-alternateur vaut \(72\%\).

Donnees : \(\rho=\SI{1000}{kg.m^{-3}}\), \(g=\SI{9.81}{m.s^{-2}}\).

\textbf{Partie A - Puissance hydraulique.}
\begin{enumerate}
\item Calculer le debit massique.
\item Justifier \(P=\rho D_vgh\).
\item Calculer la puissance hydraulique.
\end{enumerate}

\textbf{Partie B - Electricite.}
\begin{enumerate}[resume]
\item Calculer la puissance electrique produite.
\item Calculer l'energie journaliere en joules.
\item Convertir en \(\si{kWh}\).
\end{enumerate}

\textbf{Partie C - Pression.}
\begin{enumerate}[resume]
\item Calculer la surpression au bas d'une conduite statique de hauteur \(\SI{12}{m}\).
\item Pourquoi la pression en ecoulement n'est-elle pas exactement cette valeur ?
\item Quelle grandeur augmente si la section de conduite diminue ?
\end{enumerate}
\end{probleme}

\begin{probleme}{24 \diffthree}{Dissolution exothermique de soude : calorimetrie et securite}
On dissout \(\SI{8.00}{g}\) d'hydroxyde de sodium solide dans \(\SI{200}{g}\) d'eau initialement a \(\SI{20.0}{\celsius}\). La temperature finale vaut \(\SI{30.5}{\celsius}\). On assimile la solution a de l'eau.

Donnees : \(M(NaOH)=\SI{40.0}{g.mol^{-1}}\), \(c=\SI{4180}{J.kg^{-1}.K^{-1}}\).

\textbf{Partie A - Quantite.}
\begin{enumerate}
\item Calculer la quantite de \(NaOH\) dissoute.
\item Calculer la masse totale de solution.
\end{enumerate}

\textbf{Partie B - Energie.}
\begin{enumerate}[resume]
\item Calculer l'energie recue par la solution.
\item En deduire l'energie liberee par mole dissoute.
\item Qualifier la dissolution.
\end{enumerate}

\textbf{Partie C - Variante difficile.}
\begin{enumerate}[resume]
\item Si l'on dissout \(\SI{20.0}{g}\) de \(NaOH\) dans la meme masse d'eau, estimer la temperature finale en supposant la meme energie molaire.
\item Pourquoi cette extrapolation devient-elle dangereuse ?
\end{enumerate}
\end{probleme}

\subsection{Niveau \difffour{} - problemes difficiles type concours guide}

\begin{probleme}{25 \difffour}{Mission polaire : produire de l'eau liquide a partir de neige}
Une equipe polaire doit produire \(\SI{8.0}{L}\) d'eau liquide a \(\SI{35}{\celsius}\) a partir de neige assimilee a de la glace a \(\SI{-15}{\celsius}\). Le combustible libere \(\SI{42}{MJ.kg^{-1}}\), mais le rendement du rechaud n'est que \(30\%\).

Donnees : \(c_g=\SI{2100}{J.kg^{-1}.K^{-1}}\), \(L_f=\SI{334}{kJ.kg^{-1}}\), \(c_e=\SI{4180}{J.kg^{-1}.K^{-1}}\), \(\rho=\SI{1000}{kg.m^{-3}}\).

\textbf{Partie A - Modelisation.}
\begin{enumerate}
\item Convertir le volume final en masse d'eau.
\item Justifier l'approximation sur la masse de glace initiale.
\item Decomposer le processus en trois etapes thermiques.
\end{enumerate}

\textbf{Partie B - Bilan energetique.}
\begin{enumerate}[resume]
\item Calculer l'energie pour chauffer la glace a \(\SI{0}{\celsius}\).
\item Calculer l'energie de fusion.
\item Calculer l'energie pour chauffer l'eau a \(\SI{35}{\celsius}\).
\item Identifier l'etape dominante.
\end{enumerate}

\textbf{Partie C - Combustible et securite.}
\begin{enumerate}[resume]
\item Calculer l'energie chimique requise.
\item Calculer la masse minimale de combustible.
\item Proposer une marge realiste et la justifier par trois arguments.
\end{enumerate}
\end{probleme}

\begin{probleme}{26 \difffour}{Refroidissement liquide d'un serveur informatique}
Un serveur dissipe \(\SI{1.20}{kW}\). Une plaque froide fait circuler de l'eau entrant a \(\SI{18}{\celsius}\) et sortant au plus a \(\SI{32}{\celsius}\). Le tube a un diametre interieur \(\SI{6.0}{mm}\). La plaque est \(\SI{0.50}{m}\) au-dessus du reservoir.

Donnees : \(\rho=\SI{1000}{kg.m^{-3}}\), \(c=\SI{4180}{J.kg^{-1}.K^{-1}}\), \(g=\SI{9.81}{m.s^{-2}}\).

\textbf{Partie A - Debit impose par la chaleur.}
\begin{enumerate}
\item Ecrire le bilan de puissance thermique.
\item Calculer le debit massique minimal.
\item Convertir en \(\si{L.min^{-1}}\).
\end{enumerate}

\textbf{Partie B - Ecoulement.}
\begin{enumerate}[resume]
\item Calculer la section du tube.
\item Calculer la vitesse moyenne.
\item Commenter le regime attendu qualitativement.
\end{enumerate}

\textbf{Partie C - Pompe.}
\begin{enumerate}[resume]
\item Calculer la surpression hydrostatique minimale.
\item Pourquoi la pompe doit-elle fournir davantage ?
\item Que se passe-t-il si le debit est divise par deux ?
\end{enumerate}

\textbf{Partie D - Choix de l'eau.}
\begin{enumerate}[resume]
\item Pourquoi l'eau est-elle plus efficace que l'air ?
\item Citer deux risques techniques d'un refroidissement liquide.
\end{enumerate}
\end{probleme}

\begin{probleme}{27 \difffour}{Production de dihydrogene et stockage sous pression}
On produit du dihydrogene par reaction du zinc avec un acide :
\[
Zn+2H_3O^+\rightarrow Zn^{2+}+H_2+2H_2O.
\]
On introduit \(\SI{6.50}{g}\) de zinc dans \(\SI{250}{mL}\) d'acide de concentration \(\SI{1.00}{mol.L^{-1}}\). Le gaz est recueilli dans un reservoir de volume \(\SI{2.00}{L}\) a \(\SI{293}{K}\).

Donnees : \(M(Zn)=\SI{65.4}{g.mol^{-1}}\), \(R=\SI{8.314}{J.mol^{-1}.K^{-1}}\). La combustion de \(\SI{1}{mol}\) de \(H_2\) libere \(\SI{286}{kJ}\).

\textbf{Partie A - Avancement.}
\begin{enumerate}
\item Calculer les quantites initiales.
\item Identifier le reactif limitant.
\item Calculer la quantite de \(H_2\) produite.
\end{enumerate}

\textbf{Partie B - Stockage.}
\begin{enumerate}[resume]
\item Calculer la pression dans le reservoir.
\item Convertir en bar.
\item Discuter la dangerosite de la pression et du gaz.
\end{enumerate}

\textbf{Partie C - Energie.}
\begin{enumerate}[resume]
\item Calculer l'energie recuperable par combustion.
\item Comparer a l'energie pour chauffer \(\SI{1.0}{kg}\) d'eau de \(\SI{20}{\celsius}\) a \(\SI{80}{\celsius}\).
\item Pourquoi le dihydrogene est-il interessant malgre cette faible quantite ?
\end{enumerate}
\end{probleme}

\begin{probleme}{28 \difffour}{Barrage de montagne : force sur une vanne et puissance}
Un barrage alimente une conduite forcee. Le niveau de l'eau est \(\SI{85}{m}\) au-dessus de la turbine. Le debit est \(\SI{1.20}{m^3.s^{-1}}\). Une vanne verticale de largeur \(\SI{2.0}{m}\) et de hauteur \(\SI{3.0}{m}\) a son centre a \(\SI{30}{m}\) sous la surface du lac.

Donnees : \(\rho=\SI{1000}{kg.m^{-3}}\), \(g=\SI{9.81}{m.s^{-2}}\), \(p_{atm}=\SI{1.01e5}{Pa}\), rendement turbine-alternateur \(82\%\).

\textbf{Partie A - Vanne.}
\begin{enumerate}
\item Calculer la pression au centre de la vanne.
\item Calculer l'aire de la vanne.
\item Estimer la force pressante.
\item Expliquer pourquoi le calcul exact demanderait une integration.
\end{enumerate}

\textbf{Partie B - Turbine.}
\begin{enumerate}[resume]
\item Calculer le debit massique.
\item Calculer la puissance hydraulique disponible.
\item Calculer la puissance electrique.
\end{enumerate}

\textbf{Partie C - Exploitation.}
\begin{enumerate}[resume]
\item Calculer l'energie produite en une heure.
\item Convertir en \(\si{kWh}\).
\item Estimer le nombre de foyers de puissance moyenne \(\SI{1.0}{kW}\) alimentables.
\end{enumerate}
\end{probleme}

\begin{probleme}{29 \difffour}{Neutralisation concentree : chimie, temperature et risque}
On melange \(\SI{100}{mL}\) d'acide chlorhydrique \(\SI{1.00}{mol.L^{-1}}\) avec \(\SI{100}{mL}\) de soude \(\SI{1.00}{mol.L^{-1}}\), initialement a \(\SI{20.0}{\celsius}\). La neutralisation libere \(\SI{57.0}{kJ.mol^{-1}}\). On assimile la solution finale a de l'eau de masse volumique \(\SI{1000}{kg.m^{-3}}\).

\textbf{Partie A - Chimie.}
\begin{enumerate}
\item Ecrire la reaction de neutralisation.
\item Calculer les quantites initiales.
\item Determiner l'avancement maximal.
\end{enumerate}

\textbf{Partie B - Thermique.}
\begin{enumerate}[resume]
\item Calculer l'energie liberee.
\item Calculer la masse de solution finale.
\item Estimer la temperature finale theorique.
\end{enumerate}

\textbf{Partie C - Variante dangereuse.}
\begin{enumerate}[resume]
\item Refaire l'estimation si les deux solutions etaient a \(\SI{5.00}{mol.L^{-1}}\).
\item Pourquoi le modele devient-il insuffisant ?
\item Enoncer trois precautions experimentales.
\end{enumerate}
\end{probleme}

\subsection{Niveau \difffive{} - syntheses ambitieuses type concours}

\begin{probleme}{30 \difffive}{Station autonome en montagne : hydraulique, propane et bilan carbone}
Une station scientifique de montagne dispose d'une microturbine et d'un bruleur au propane. L'eau alimente la turbine avec un debit \(D_v=\SI{0.060}{m^3.s^{-1}}\) et une chute de \(\SI{18}{m}\). Le rendement turbine-alternateur est \(65\%\). Le bruleur chauffe \(\SI{40}{L}\) d'eau de \(\SI{5}{\celsius}\) a \(\SI{60}{\celsius}\). Le propane libere \(\SI{2220}{kJ.mol^{-1}}\), rendement thermique \(50\%\).

Donnees : \(\rho=\SI{1000}{kg.m^{-3}}\), \(c=\SI{4180}{J.kg^{-1}.K^{-1}}\), \(g=\SI{9.81}{m.s^{-2}}\), \(M(C_3H_8)=\SI{44.0}{g.mol^{-1}}\), \(M(CO_2)=\SI{44.0}{g.mol^{-1}}\).

\textbf{Partie A - Microturbine.}
\begin{enumerate}
\item Calculer le debit massique.
\item Calculer la puissance hydraulique.
\item Calculer la puissance electrique produite.
\item Calculer l'energie electrique produite en \(\SI{2.0}{h}\).
\end{enumerate}

\textbf{Partie B - Chauffage au propane.}
\begin{enumerate}[resume]
\item Calculer l'energie utile pour chauffer l'eau.
\item Calculer l'energie chimique a fournir.
\item Calculer la masse de propane consommee.
\end{enumerate}

\textbf{Partie C - Bilan carbone.}
\begin{enumerate}[resume]
\item Ajuster la combustion du propane.
\item Calculer la masse de \(CO_2\) rejetee.
\item Comparer qualitativement l'usage turbine et l'usage propane.
\end{enumerate}

\textbf{Partie D - Strategie energetique.}
\begin{enumerate}[resume]
\item Comparer energie electrique produite et energie thermique utile.
\item Pourquoi ces energies ne sont-elles pas directement equivalentes en pratique ?
\item Proposer trois ameliorations techniques.
\end{enumerate}
\end{probleme}

\begin{probleme}{31 \difffive}{Fusee a eau : air comprime, debit initial et poussee}
Une fusee a eau est une bouteille de volume total \(\SI{1.50}{L}\), contenant initialement \(\SI{0.60}{L}\) d'eau et de l'air comprime a \(\SI{5.0}{bar}\) absolus, a \(\SI{293}{K}\). L'eau sort par une tuyere de diametre \(\SI{8.0}{mm}\). On suppose initialement :
\[
\Delta p \simeq \frac12\rho v^2.
\]
La masse totale de la fusee au depart est \(\SI{0.90}{kg}\).

Donnees : \(p_{atm}=\SI{1.0}{bar}\), \(\rho=\SI{1000}{kg.m^{-3}}\), \(R=\SI{8.314}{J.mol^{-1}.K^{-1}}\), \(g=\SI{9.81}{m.s^{-2}}\).

\textbf{Partie A - Air comprime.}
\begin{enumerate}
\item Calculer le volume initial d'air.
\item Calculer la quantite d'air comprime.
\item Comparer a la quantite d'air a pression atmospherique dans le meme volume.
\end{enumerate}

\textbf{Partie B - Jet initial.}
\begin{enumerate}[resume]
\item Calculer la surpression initiale.
\item Estimer la vitesse initiale du jet.
\item Calculer la section de la tuyere.
\item Estimer le debit volumique initial.
\end{enumerate}

\textbf{Partie C - Poussee.}
\begin{enumerate}[resume]
\item Calculer le debit massique initial.
\item Estimer la poussee \(F\simeq D_m v\).
\item Comparer au poids de la fusee.
\end{enumerate}

\textbf{Partie D - Critique.}
\begin{enumerate}[resume]
\item Pourquoi la poussee diminue-t-elle ?
\item Pourquoi l'hypothese de pression constante est-elle fausse ?
\item Citer trois parametres influencant la hauteur atteinte.
\end{enumerate}
\end{probleme}

\begin{probleme}{32 \difffive}{Hopital de campagne : chauffage, perfusion et dilution}
Un hopital de campagne doit chauffer \(\SI{25}{L}\) d'eau de \(\SI{15}{\celsius}\) a \(\SI{90}{\celsius}\), installer une perfusion gravitaire demandant une surpression de \(\SI{2.8e3}{Pa}\), et preparer \(\SI{2.00}{L}\) d'une solution desinfectante de concentration \(\SI{2.0e-2}{mol.L^{-1}}\) a partir d'une solution mere de concentration \(\SI{0.80}{mol.L^{-1}}\). Le bruleur utilise de l'ethanol, dont la combustion libere \(\SI{1367}{kJ.mol^{-1}}\), avec un rendement \(40\%\).

Donnees : \(c=\SI{4180}{J.kg^{-1}.K^{-1}}\), \(\rho_{eau}=\SI{1000}{kg.m^{-3}}\), \(\rho_{perf}=\SI{1030}{kg.m^{-3}}\), \(M(C_2H_6O)=\SI{46.0}{g.mol^{-1}}\), \(M(CO_2)=\SI{44.0}{g.mol^{-1}}\).

\textbf{Partie A - Chauffage.}
\begin{enumerate}
\item Calculer la masse d'eau.
\item Calculer l'energie utile.
\item Calculer la masse d'ethanol necessaire.
\end{enumerate}

\textbf{Partie B - Gaz de combustion.}
\begin{enumerate}[resume]
\item Ajuster la combustion de l'ethanol.
\item Calculer la masse de \(CO_2\) produite.
\item Pourquoi faut-il ventiler ?
\end{enumerate}

\textbf{Partie C - Perfusion.}
\begin{enumerate}[resume]
\item Calculer la hauteur minimale de la poche.
\item Expliquer pourquoi le debit depend aussi du catheter.
\end{enumerate}

\textbf{Partie D - Dilution.}
\begin{enumerate}[resume]
\item Calculer la quantite de solute finale.
\item Calculer le volume de solution mere a prelever.
\item Rediger le protocole de dilution.
\end{enumerate}

\textbf{Partie E - Synthese.}
\begin{enumerate}[resume]
\item Associer chaque tache au chapitre principal mobilise.
\item Expliquer pourquoi un probleme reel melange souvent plusieurs domaines.
\end{enumerate}
\end{probleme}

\newpage

% ============================================================
% CORRIGES
% ============================================================

\section{Corriges detailles}

\begin{corrige}
\textbf{Probleme 1.} 
\[
V_m=0.250\,L,\quad n=C_mV_m=0.100\times0.250=2.50\times10^{-2}\,mol.
\]
\[
m=nM=2.50\times10^{-2}\times58.5=1.46\,g.
\]
Protocole : peser \(\SI{1.46}{g}\), introduire dans une fiole jaugee de \(\SI{250}{mL}\), dissoudre, completer au trait, homogeniser. Le facteur de dilution vaut \(F=C_m/C_f=10\), donc il faut prelever \(V_m'=V_f/F=\SI{10.0}{mL}\). Une pipette jaugee limite l'incertitude sur le volume preleve. Pour la conductivite : \(C_f=10.0\,mol.m^{-3}\), donc
\[
\sigma_{th}=(5.01+7.63)10^{-3}\times10.0=0.126\,S.m^{-1}.
\]
Ecart relatif :
\[
\frac{|0.119-0.126|}{0.126}=5.6\%.
\]
Causes possibles : temperature, dilution imparfaite, etalonnage du conductimetre, purete du sel.
\end{corrige}

\begin{corrige}
\textbf{Probleme 2.}
\[
Q=mc\Delta T=0.750\times4180\times(95-18)=2.41\times10^5\,J.
\]
\[
E_{elec}=P\Delta t=1800\times160=2.88\times10^5\,J.
\]
\[
\eta=Q/E_{elec}=0.837\simeq84\%.
\]
Pour la seconde chauffe :
\[
Q'=1.20\times4180\times80=4.01\times10^5\,J,
\quad E'=Q'/\eta=4.79\times10^5\,J.
\]
\[
t=E'/P=266\,s\simeq4.4\,min.
\]
Le rendement varie car les pertes dependent de la duree, du recipient et de l'ecart avec l'air. L'ebullition depend aussi de la pression atmospherique.
\end{corrige}

\begin{corrige}
\textbf{Probleme 3.}
\[
\Delta p=\rho gh=1000\times9.81\times3.20=3.14\times10^4\,Pa.
\]
\[
p=1.01\times10^5+3.14\times10^4=1.32\times10^5\,Pa=1.32\,bar.
\]
\[
S=\pi(0.120)^2=4.52\times10^{-2}\,m^2.
\]
Force absolue : \(F=pS=5.97\times10^3\,N\). Force due a la seule surpression :
\[
F_{sur}=\Delta p\,S=3.14\times10^4\times4.52\times10^{-2}=1.42\times10^3\,N.
\]
Le poids d'une personne de \(70\,kg\) vaut \(687\,N\), donc la surpression equivaut a plus de deux personnes. La pression depend de la hauteur de colonne d'eau, pas de la largeur du bassin.
\end{corrige}

\begin{corrige}
\textbf{Probleme 4.}
\[
H_3O^+ + HO^-\rightarrow2H_2O.
\]
A l'equivalence : \(n(H_3O^+)=n(HO^-)\).
\[
n=C_BV_E=0.100\times14.8\times10^{-3}=1.48\times10^{-3}\,mol.
\]
Dans \(10.0\,mL\) dilues :
\[
C_{dil}=\frac{1.48\times10^{-3}}{10.0\times10^{-3}}=0.148\,mol.L^{-1}.
\]
Facteur de dilution \(F=100.0/5.00=20.0\), donc \(C=2.96\,mol.L^{-1}\). Concentration massique :
\[
C_m=2.96\times36.5=108\,g.L^{-1}.
\]
L'annonce \(110\,g.L^{-1}\) est coherente. Precautions : lunettes, gants, dilution lente, pipetage avec propipette.
\end{corrige}

\begin{corrige}
\textbf{Probleme 5.}
\[
D_v=1.50/12.0=0.125\,L.s^{-1}=1.25\times10^{-4}\,m^3.s^{-1}.
\]
\[
S=\pi(4.0\times10^{-3})^2=5.03\times10^{-5}\,m^2,
\quad v=D_v/S=2.49\,m.s^{-1}.
\]
En trois heures :
\[
V=0.125\times3.0\times3600=1350\,L=1.35\,m^3.
\]
En trente jours : \(40.5\,m^3\). C'est surestime si la fontaine n'est pas ouverte en permanence ou si le debit varie.
\end{corrige}

\begin{corrige}
\textbf{Probleme 6.}
\[
Q_f=0.040\times334000=1.34\times10^4\,J.
\]
\[
Q_{eau,max}=0.300\times4180\times22=2.76\times10^4\,J.
\]
Toute la glace fond. Il reste \(1.42\times10^4\,J\) pour rechauffer \(0.340\,kg\) d'eau liquide :
\[
T_f=\frac{1.42\times10^4}{0.340\times4180}=10.0^\circ C.
\]
Si la glace est doublee, fusion necessaire \(2.67\times10^4\,J\), presque toute l'energie de l'eau sert a fondre. Il reste environ \(27588-26720=868\,J\), donc temperature finale proche de \(0.7^\circ C\), sans glace restante notable dans ce modele.
\end{corrige}

\begin{corrige}
\textbf{Probleme 7.}
\[
F_A=\rho_{air}gV=1.20\times9.81\times0.030=0.353\,N.
\]
\[
m_{He}=0.18\times0.030=5.4\times10^{-3}\,kg.
\]
Masse sans charge : \(0.0080+0.0054=0.0134\,kg\), poids \(0.131\,N\). Force disponible : \(0.222\,N\), donc masse maximale :
\[
m=0.222/9.81=2.26\times10^{-2}\,kg=22.6\,g.
\]
Cette valeur est ideale : ficelle, variations de densite, fuites et courants d'air diminuent la charge utile.
\end{corrige}

\begin{corrige}
\textbf{Probleme 8.}
\[
C_3H_8+5O_2\rightarrow3CO_2+4H_2O.
\]
\[
n(C_3H_8)=5.00/44.0=0.1136\,mol.
\]
\[
n(CO_2)=3n=0.341\,mol,\quad m(CO_2)=0.341\times44.0=15.0\,g.
\]
\[
E=0.1136\times2220=252\,kJ,\quad E_u=0.45E=113\,kJ.
\]
\[
m_{eau}=\frac{113000}{4180\times60}=0.451\,kg.
\]
Le rendement est limite par les pertes vers l'air, la casserole et les gaz chauds. En tente fermee, risques de \(CO_2\), manque de \(O_2\), voire \(CO\) si combustion incomplete.
\end{corrige}

\begin{corrige}
\textbf{Probleme 9.}
A l'equivalence : \(n_A=n_B\). Donc
\[
n_B=0.100\times16.4\times10^{-3}=1.64\times10^{-3}\,mol.
\]
Dans \(20.0\,mL\) :
\[
C_{dil}=\frac{1.64\times10^{-3}}{20.0\times10^{-3}}=8.20\times10^{-2}\,mol.L^{-1}.
\]
Facteur \(10\), donc \(C_{vin}=0.820\,mol.L^{-1}\). Concentration massique :
\[
C_m=0.820\times60.0=49.2\,g.L^{-1}.
\]
Ecart avec \(50\,g.L^{-1}\) : \(1.6\%\). Sources : burette, pipette, equivalence, concentration de soude. La dilution rend le volume equivalent plus confortable et limite les erreurs relatives.
\end{corrige}

\begin{corrige}
\textbf{Probleme 10.}
\[
m=\rho V=1000\times0.150=150\,kg.
\]
\[
Q=150\times4180\times34=2.13\times10^7\,J=21.3\,MJ=5.92\,kWh.
\]
\[
\eta=21.3/26=0.819.
\]
Pour la douche, si \(V_h\) est le volume chaud :
\[
V_h(52-40)=(45-V_h)(40-15).
\]
\[
12V_h=1125-25V_h,\quad V_h=30.4\,L.
\]
Nombre de douches : \(150/30.4=4.9\), donc quatre a cinq selon marge.
\end{corrige}

\begin{corrige}
\textbf{Probleme 11.}
\[
h=\frac{\Delta p}{\rho g}=\frac{3.0\times10^3}{1030\times9.81}=0.297\,m.
\]
Debit :
\[
D_v=\frac{2.5\times10^{-6}}{60}=4.17\times10^{-8}\,m^3.s^{-1}.
\]
\[
v=D_v/S=\frac{4.17\times10^{-8}}{0.75\times10^{-6}}=5.56\times10^{-2}\,m.s^{-1}.
\]
Pour \(500\,mL\) : \(t=500/2.5=200\,min\). La hauteur ne suffit pas car il existe viscosite, frottements, pertes de charge, catheter et pression veineuse. Si le bras monte, la difference de hauteur diminue, donc la surpression aussi.
\end{corrige}

\begin{corrige}
\textbf{Probleme 12.}
Masse de carbonate : \(0.75\times1.80=1.35\,g\). Donc
\[
n(CaCO_3)=1.35/100=1.35\times10^{-2}\,mol.
\]
\[
n(H_3O^+)=0.500\times0.100=5.00\times10^{-2}\,mol.
\]
Avancements :
\[
x_{carbonate}=1.35\times10^{-2},\quad x_{acide}=\frac{5.00\times10^{-2}}{2}=2.50\times10^{-2}.
\]
Le carbonate est limitant. Donc \(n(CO_2)=1.35\times10^{-2}\,mol\), \(m=0.594\,g\), et
\[
V=1.35\times10^{-2}\times24.0=0.324\,L.
\]
L'effervescence est due au degagement du gaz. Avec un acide deux fois moins concentre, l'acide resterait encore en exces legerement : \(x_{acide}=1.25\times10^{-2}\), il deviendrait limitant.
\end{corrige}

\begin{corrige}
\textbf{Probleme 13.}
\[
\dot m=\frac{P}{c\Delta T}=\frac{45}{1000\times16}=2.81\times10^{-3}\,kg.s^{-1}.
\]
\[
D_v=\dot m/\rho=2.34\times10^{-3}\,m^3.s^{-1}.
\]
\[
S=12\,cm^2=1.2\times10^{-3}\,m^2,
\quad v=D_v/S=1.95\,m.s^{-1}.
\]
La vitesse est plausible. La temperature de sortie varie selon les zones du radiateur. Conduction dans les composants, rayonnement et convection naturelle sont oublies. La poussiere diminue la section efficace et isole thermiquement.
\end{corrige}

\begin{corrige}
\textbf{Probleme 14.}
\[
V=1.20\times10^{-2}\,m^3,\quad p=2.00\times10^7\,Pa.
\]
\[
n=\frac{pV}{RT}=\frac{2.00\times10^7\times1.20\times10^{-2}}{8.314\times293}=98.6\,mol.
\]
Volume equivalent : \(200\times12.0=2400\,L\). Autonomie surface : \(2400/20=120\,min\). A \(20\,m\), pression absolue environ \(3\,bar\), donc la consommation en quantite de matiere est environ triple. On garde une reserve de securite et on remonte lentement pour eviter les risques lies a la decompression.
\end{corrige}

\begin{corrige}
\textbf{Probleme 15.}
\[
H_3O^+ + HO^-\rightarrow2H_2O.
\]
\[
n=CV=2.0\times10^{-2}\times80=1.60\,mol.
\]
\[
V_B=n/C_B=1.60/0.50=3.20\,L.
\]
Energie liberee : \(1.60\times57=91.2\,kJ\). Elevation :
\[
\Delta T=\frac{91200}{80\times4180}=0.273\,K.
\]
Un exces de soude rend le rejet basique. On verifie le pH final. Precautions : ajout lent, agitation, lunettes, gants.
\end{corrige}

\begin{corrige}
\textbf{Probleme 16.}
\[
\Delta p=\rho gh=1000\times9.81\times6.0=5.89\times10^4\,Pa=0.589\,bar.
\]
Sous hypotheses de Bernoulli : fluide parfait, reservoir grand, pertes negligees, sortie a l'air libre.
\[
v=\sqrt{2gh}=\sqrt{2\times9.81\times6.0}=10.8\,m.s^{-1}.
\]
\[
S=\pi(5.0\times10^{-3})^2=7.85\times10^{-5}\,m^2.
\]
\[
D_v=Sv=8.48\times10^{-4}\,m^3.s^{-1}=50.9\,L.min^{-1}.
\]
Le debit reel est diminue par les pertes. Un tuyau long et un embout fin augmentent les pertes.
\end{corrige}

\begin{corrige}
\textbf{Probleme 17.}
\[
n=C_fV=2.0\times10^{-4}\times500=0.100\,mol.
\]
\[
V_0=n/C_0=0.100/0.50=0.200\,L.
\]
Solution dix fois diluee : \(0.050\,mol.L^{-1}\), volume a ajouter : \(0.100/0.050=2.00\,L\). Cette mesure est plus robuste. Controle :
\[
C=\frac{9.0\times10^{-6}}{50.0\times10^{-3}}=1.8\times10^{-4}\,mol.L^{-1}.
\]
Ecart relatif : \(10\%\). Causes : decomposition, reaction avec matiere organique, melange imparfait, dosage.
\end{corrige}

\begin{corrige}
\textbf{Probleme 18.}
\[
Q_{eau}=2.00\times4180\times40=3.34\times10^5\,J.
\]
\[
E=600\times12.0\times60=4.32\times10^5\,J,
\quad \eta=0.774.
\]
Recipient :
\[
Q_{rec}=650\times40=2.60\times10^4\,J.
\]
\[
\eta'=(3.34\times10^5+2.60\times10^4)/(4.32\times10^5)=0.833.
\]
La puissance de maintien compense les pertes. En \(30\,min\), \(E=90\times1800=1.62\times10^5\,J\). Les pertes augmentent avec l'ecart de temperature avec l'air.
\end{corrige}

\begin{corrige}
\textbf{Probleme 19.}
\[
n_0=1500/180=8.33\,mol,
\quad n_c=0.65\times8.33=5.42\,mol.
\]
\[
n(CO_2)=2n_c=10.8\,mol.
\]
\[
m(CO_2)=10.8\times44.0=475\,g,
\quad V=10.8\times24.0=260\,L.
\]
Il faut evacuer le gaz pour eviter une surpression. Energie thermique :
\[
E=5.42\times70=379\,kJ.
\]
\[
\Delta T=379000/(20\times4000)=4.74\,K.
\]
Un controle thermique preserve les levures et la qualite du produit.
\end{corrige}

\begin{corrige}
\textbf{Probleme 20.}
\[
D_v=6.0\times10^{-3}/60=1.00\times10^{-4}\,m^3.s^{-1}.
\]
\[
\dot m=\rho D_v=0.100\,kg.s^{-1}.
\]
\[
P_u=\dot m c\Delta T=0.100\times4180\times30=12.5\,kW.
\]
\[
P_e=12.5/0.92=13.6\,kW.
\]
Une prise de \(3.7\,kW\) ne suffit pas. A \(3.7\,kW\), puissance utile \(3.40\,kW\), donc
\[
\dot m=\frac{3400}{4180\times30}=2.71\times10^{-2}\,kg.s^{-1},
\]
soit \(1.63\,L.min^{-1}\). Une installation specifique est necessaire pour les fortes puissances.
\end{corrige}

\begin{corrige}
\textbf{Probleme 21.}
\[
P=70.0\times9.81=687\,N.
\]
\[
F_A=1000\times9.81\times0.080=785\,N.
\]
Le sous-marin remonte. Pour l'equilibre : masse totale \(=\rho V=80.0\,kg\), donc \(m_b=10.0\,kg\), soit \(10.0\,L\). Le ballast maximal suffit. A \(25\,m\) :
\[
p=1.01\times10^5+1000\times9.81\times25=3.46\times10^5\,Pa.
\]
\[
S=\pi(0.080)^2=2.01\times10^{-2}\,m^2,
\quad F=pS=6.95\times10^3\,N.
\]
La force est de plusieurs kilonewtons : le hublot doit etre resistant.
\end{corrige}

\begin{corrige}
\textbf{Probleme 22.}
Combustion :
\[
2C_8H_{18}+25O_2\rightarrow16CO_2+18H_2O.
\]
\[
m=750\times0.80\times10^{-3}=0.600\,kg.
\]
\[
n=600/114=5.26\,mol.
\]
\[
E=5.26\times5470=28.8\,MJ,
\quad P_{chim}=28.8\times10^6/3600=8.00\,kW.
\]
\[
P_m=2.00\,kW,
\quad P_{th}=6.00\,kW.
\]
Debit d'eau :
\[
\dot m=\frac{6000}{4180\times50}=2.87\times10^{-2}\,kg.s^{-1}=1.72\,L.min^{-1}.
\]
Le radiateur evacue vers l'air, le thermostat stabilise la temperature.
\end{corrige}

\begin{corrige}
\textbf{Probleme 23.}
\[
\dot m=\rho D_v=800\,kg.s^{-1}.
\]
\[
P=\rho D_vgh=1000\times0.80\times9.81\times12=9.42\times10^4\,W.
\]
\[
P_e=0.72P=6.78\times10^4\,W.
\]
Energie jour :
\[
E=6.78\times10^4\times86400=5.86\times10^9\,J=1627\,kWh.
\]
Surpression statique :
\[
\Delta p=1000\times9.81\times12=1.18\times10^5\,Pa.
\]
En ecoulement, energie cinetique et pertes modifient la pression. Si la section diminue, la vitesse augmente.
\end{corrige}

\begin{corrige}
\textbf{Probleme 24.}
\[
n=8.00/40.0=0.200\,mol,
\quad m_{sol}=0.208\,kg.
\]
\[
Q=0.208\times4180\times10.5=9.13\times10^3\,J.
\]
\[
Q_m=9.13/0.200=45.6\,kJ.mol^{-1}.
\]
La dissolution est exothermique. Pour \(20.0\,g\), \(n=0.500\,mol\), energie \(22.8\,kJ\). Masse solution \(0.220\,kg\), donc
\[
\Delta T=22800/(0.220\times4180)=24.8\,K,
\]
soit \(44.8^\circ C\). L'extrapolation est dangereuse car les capacites changent, les pertes aussi, et la soude concentree est corrosive.
\end{corrige}

\begin{corrige}
\textbf{Probleme 25.}
Masse d'eau : \(8.0\,kg\). Etapes : chauffer glace, fondre, chauffer eau.
\[
Q_1=8.0\times2100\times15=2.52\times10^5\,J.
\]
\[
Q_2=8.0\times334000=2.67\times10^6\,J.
\]
\[
Q_3=8.0\times4180\times35=1.17\times10^6\,J.
\]
Total \(Q=4.09\,MJ\), etape dominante : fusion. Energie chimique :
\[
E=Q/0.30=13.6\,MJ.
\]
\[
m=13.6/42=0.325\,kg.
\]
Prevoir plus : vent, pertes, casserole froide, neige plus froide, rendement variable.
\end{corrige}

\begin{corrige}
\textbf{Probleme 26.}
\[
\dot m=\frac{1200}{4180\times14}=2.05\times10^{-2}\,kg.s^{-1}.
\]
\[
D_v=2.05\times10^{-5}\,m^3.s^{-1}=1.23\,L.min^{-1}.
\]
\[
S=\pi(3.0\times10^{-3})^2=2.83\times10^{-5}\,m^2,
\quad v=0.724\,m.s^{-1}.
\]
\[
\Delta p=1000\times9.81\times0.50=4.91\times10^3\,Pa.
\]
La pompe doit vaincre pertes de charge et obstacles. Si debit divise par deux, \(\Delta T\) double : sortie vers \(46^\circ C\). L'eau est efficace grace a son grand \(c\), mais risques de fuite, corrosion, encrassement.
\end{corrige}

\begin{corrige}
\textbf{Probleme 27.}
\[
n(Zn)=6.50/65.4=9.94\times10^{-2}\,mol.
\]
\[
n(H_3O^+)=1.00\times0.250=0.250\,mol.
\]
\[
x_{Zn}=9.94\times10^{-2},\quad x_{acide}=0.250/2=0.125.
\]
Zinc limitant, donc \(n(H_2)=9.94\times10^{-2}\,mol\). Pression :
\[
p=\frac{nRT}{V}=\frac{9.94\times10^{-2}\times8.314\times293}{2.00\times10^{-3}}=1.21\times10^5\,Pa=1.21\,bar.
\]
Energie combustion :
\[
E=0.0994\times286=28.4\,kJ.
\]
Chauffer \(1\,kg\) d'eau de \(60\,K\) demande \(251\,kJ\), donc ici c'est insuffisant. L'interet du dihydrogene tient a sa densite energetique massique et a son usage en pile, mais son stockage est delicat.
\end{corrige}

\begin{corrige}
\textbf{Probleme 28.}
\[
p=1.01\times10^5+1000\times9.81\times30=3.95\times10^5\,Pa.
\]
\[
S=2.0\times3.0=6.0\,m^2,
\quad F=pS=2.37\times10^6\,N.
\]
Le calcul exact integre une pression variable avec la profondeur. Debit massique : \(1200\,kg.s^{-1}\). Puissance hydraulique :
\[
P=1200\times9.81\times85=1.00\times10^6\,W.
\]
Puissance electrique : \(0.82\,MW\). En une heure : \(820\,kWh\), soit environ \(820\) foyers de \(1.0\,kW\) moyen.
\end{corrige}

\begin{corrige}
\textbf{Probleme 29.}
\[
n(H_3O^+)=n(HO^-)=1.00\times0.100=0.100\,mol.
\]
\[
E=0.100\times57.0=5.70\,kJ.
\]
Masse finale \(0.200\,kg\). Donc
\[
\Delta T=5700/(0.200\times4180)=6.82\,K,
\quad T_f=26.8^\circ C.
\]
A \(5.00\,mol.L^{-1}\), \(n=0.500\,mol\), energie \(28.5\,kJ\), donc \(\Delta T=34.1\,K\), \(T_f\approx54^\circ C\). Modele insuffisant : pertes, variation de capacite, projections, melange non instantane. Precautions : lunettes, gants, ajout lent, agitation, bain froid si besoin.
\end{corrige}

\begin{corrige}
\textbf{Probleme 30.}
Microturbine :
\[
\dot m=1000\times0.060=60\,kg.s^{-1}.
\]
\[
P_h=60\times9.81\times18=1.06\times10^4\,W,
\quad P_e=0.65P_h=6.89\,kW.
\]
En \(2.0\,h\) : \(13.8\,kWh\). Chauffage :
\[
Q=40\times4180\times55=9.20\,MJ.
\]
Energie chimique propane : \(18.4\,MJ\), donc
\[
n=18.4/2.220=8.29\,mol,
\quad m=8.29\times44.0=365\,g.
\]
Combustion :
\[
C_3H_8+5O_2\rightarrow3CO_2+4H_2O.
\]
\[
n(CO_2)=24.9\,mol,
\quad m(CO_2)=1.10\,kg.
\]
L'hydraulique produit une energie electrique sans combustion locale ; le propane fournit directement de la chaleur mais emet du \(CO_2\). Les energies ne sont pas equivalentes car stockage, puissance, rendement et usages different. Ameliorations : isolation, recuperation de chaleur, meilleur rendement, stockage batterie.
\end{corrige}

\begin{corrige}
\textbf{Probleme 31.}
Volume d'air : \(0.90\,L=9.0\times10^{-4}\,m^3\). Quantite d'air comprime :
\[
n=\frac{5.0\times10^5\times9.0\times10^{-4}}{8.314\times293}=0.185\,mol.
\]
A pression atmospherique : \(0.0370\,mol\), donc cinq fois moins. Surpression : \(4.0\times10^5\,Pa\). Vitesse :
\[
v=\sqrt{\frac{2\Delta p}{\rho}}=\sqrt{800}=28.3\,m.s^{-1}.
\]
\[
S=\pi(4.0\times10^{-3})^2=5.03\times10^{-5}\,m^2,
\quad D_v=1.42\times10^{-3}\,m^3.s^{-1}.
\]
\[
D_m=1.42\,kg.s^{-1},
\quad F\simeq D_mv=40.2\,N.
\]
Poids : \(0.90\times9.81=8.83\,N\). La poussee initiale est largement superieure au poids, mais elle diminue car la pression baisse et l'eau s'echappe. Parametres importants : pression initiale, volume d'eau, diametre de tuyere, masse, aerodynamique.
\end{corrige}

\begin{corrige}
\textbf{Probleme 32.}
Chauffage : masse \(25\,kg\), \(\Delta T=75\,K\).
\[
Q=25\times4180\times75=7.84\,MJ.
\]
Energie chimique : \(Q/0.40=19.6\,MJ\). Ethanol :
\[
n=19.6/1.367=14.3\,mol,
\quad m=14.3\times46.0=658\,g.
\]
Combustion :
\[
C_2H_6O+3O_2\rightarrow2CO_2+3H_2O.
\]
\[
n(CO_2)=28.6\,mol,
\quad m=28.6\times44.0=1.26\,kg.
\]
Ventilation necessaire : consommation de \(O_2\), rejet de \(CO_2\), risque de \(CO\). Perfusion :
\[
h=\frac{2.8\times10^3}{1030\times9.81}=0.277\,m.
\]
Le debit depend du catheter, des pertes, de la viscosite et de la pression veineuse. Dilution :
\[
n_f=2.0\times10^{-2}\times2.00=4.0\times10^{-2}\,mol,
\quad V_0=n_f/C_0=0.050\,L=50\,mL.
\]
Protocole : prelever \(50\,mL\), introduire dans fiole jaugee de \(2.00\,L\), completer, homogeniser. Chauffage : energie ; perfusion : hydrostatique ; dilution : conservation de matiere. Les problemes reels melangent les domaines car les objets techniques echangent matiere, energie et forces.
\end{corrige}

\newpage

% ============================================================
% FICHE METHODE
% ============================================================

\section{Fiche methode finale}

\begin{retenir}
\textbf{Chimie - reflexes.}
\begin{enumerate}
\item Equation ajustee obligatoire.
\item Convertir masses et volumes en quantites de matiere.
\item Utiliser les coefficients stoechiometriques.
\item En titrage : raisonner a l'equivalence.
\item En dilution : conserver la quantite de matiere du solute.
\end{enumerate}
\end{retenir}

\begin{retenir}
\textbf{Thermodynamique - reflexes.}
\begin{enumerate}
\item Distinguer chauffage et changement d'etat.
\item Utiliser \(Q=mc\Delta T\) hors changement d'etat.
\item Utiliser \(Q=mL\) pour fusion ou vaporisation.
\item Integrer le rendement seulement entre energie fournie et energie utile.
\item Dans un melange isole : l'energie cedee par les corps chauds est recue par les corps froids.
\end{enumerate}
\end{retenir}

\begin{retenir}
\textbf{Fluides - reflexes.}
\begin{enumerate}
\item Hydrostatique : \(p=p_0+\rho gh\).
\item Force pressante : \(F=pS\).
\item Flottabilite : comparer poids et poussee d'Archimede.
\item Debit : \(D_v=Sv\).
\item Bernoulli : seulement si les hypotheses sont explicitement acceptees.
\end{enumerate}
\end{retenir}

\begin{vigilance}
\textbf{Strategie de concours.} Dans une question longue, la difficulte n'est pas la formule mais l'enchainement. Ecrire les sous-resultats intermediaires avec leurs unites : \(n\), \(Q\), \(D_v\), \(\dot m\), \(p\), \(F\). Ces grandeurs servent souvent dans la partie suivante.
\end{vigilance}

\begin{erreur}
\textbf{Erreurs a bannir.}
\begin{itemize}
\item Grammes dans \(mc\Delta T\) avec \(c\) en \(\si{J.kg^{-1}.K^{-1}}\).
\item Litres dans une formule de pression ou de debit SI.
\item Oubli de la pression atmospherique quand on demande une pression absolue.
\item Oubli du facteur de dilution.
\item Oubli du coefficient 2 ou 3 dans une equation chimique.
\item Conclusion sans unite ni phrase physique.
\end{itemize}
\end{erreur}

\end{document}
